% GENERATED FILE. Edit canonical lessons, then run npm run generate:books.
% canonical-source-hash: fnv1a64:3278033154e98e4e
% canonical-lessons: AR-C284-awwalan-thumma, AR-C284-haka, AR-C284-indama, AR-C284-qissat-rihla, AR-R284-a-story-in-order

\chapter{A Story, In Order}
\label{ch:ar-a-story-in-order}
\hlchaptermodality{284}

\begin{chapteropening}
\textbf{By the end of this chapter:} \emph{I can tell a short story about something that happened, in order, and read one of about eighty words.}

Complete AR-R284-a-story-in-order: a short story told aloud in order, with awwalan, thumma, ʿindamā and fī an-nihāya.
\end{chapteropening}

\section[\texorpdfstring{\ar{أولا،} \ar{ثم،} \ar{بعد} \ar{ذلك} (awwalan, thumma …)}{أولا، ثم، بعد ذلك (awwalan, thumma …)}]{\ar{أولا،} \ar{ثم،} \ar{بعد} \ar{ذلك} (awwalan, thumma, baʿda dhālika) — first, then, after that}

\label{lesson:AR-C284-awwalan-thumma}



\begin{quote}
Say \textquotedblleft{}we went\textquotedblright{}, \textquotedblleft{}we ate\textquotedblright{}, \textquotedblleft{}we came back\textquotedblright{}. (\emph{Dhahabnā}, \emph{akalnā}, \emph{rajaʿnā}.)

Three events in a row are a list. A few small words turn them into a story, and you already own most of them.
\end{quote}

\begin{grammarlens}[title={Grammar Lens: first, then, after that}]
\emph{Awwalan} opens the first step. \emph{Thumma} moves to the next, with a little time between. \emph{Baʿda dhālika}, \textquotedblleft{}after that\textquotedblright{}, is \emph{baʿda} with \emph{dhālika}, \textquotedblleft{}that\textquotedblright{}, the same \emph{dhālika} that sits inside \emph{li-dhālika}.

\begin{quote}
\textbf{\ar{أولا} \ar{ذهبنا} \ar{إلى} \ar{السوق،} \ar{ثم} \ar{أكلنا} \ar{في} \ar{مطعم،} \ar{وبعد} \ar{ذلك} \ar{رجعنا} \ar{إلى} \ar{الفندق}.} — \emph{awwalan dhahabnā ilā as-sūq, thumma akalnā fī maṭʿam, wa-baʿda dhālika rajaʿnā ilā al-funduq} — First we went to the market, then we ate in a restaurant, and after that we went back to the hotel.
\end{quote}

Each joiner stands at the head of its own step, and each step has its own past verb.
\end{grammarlens}

\begin{grammarlens}[title={Grammar Lens: closing the story}]
\textbf{\ar{في} \ar{النهاية}} — \emph{fī an-nihāya} — \textquotedblleft{}in the end\textquotedblright{}. \emph{Nihāya} is \textquotedblleft{}an end\textquotedblright{}, from the root of \emph{intahā}, \textquotedblleft{}to end\textquotedblright{}.

\begin{quote}
\textbf{\ar{في} \ar{النهاية} \ar{وجدت} \ar{المفتاح}.} — \emph{fī an-nihāya wajadtu al-miftāḥ} — In the end I found the key.
\end{quote}

\emph{Akhīran}, \textquotedblleft{}at last\textquotedblright{}, closes a story too, and it usually carries relief: the thing you were waiting for came.
\end{grammarlens}

\subsection*{Your turn}
Say these aloud:

\begin{itemize}
\raggedright
  \item \emph{awwalan … thumma … baʿda dhālika … fī an-nihāya}
  \item \textquotedblleft{}First we went to the market\textquotedblright{} — \emph{awwalan dhahabnā ilā as-sūq}
  \item \textquotedblleft{}then we ate in a restaurant\textquotedblright{} — \emph{thumma akalnā fī maṭʿam}
  \item \textquotedblleft{}and after that we went back to the hotel\textquotedblright{} — \emph{wa-baʿda dhālika rajaʿnā ilā al-funduq}
  \item \textquotedblleft{}In the end I found the key\textquotedblright{} — \emph{fī an-nihāya wajadtu al-miftāḥ}
\end{itemize}

\begin{tcolorbox}[breakable,colback=teal!4,colframe=teal!35!black,title={Before you move on}]
\textquotedblleft{}After that\textquotedblright{}? (\emph{\textbf{Baʿda dhālika.}}) \textquotedblleft{}In the end\textquotedblright{}? (\emph{\textbf{Fī an-nihāya.}}) Which root does \emph{nihāya} share with \textquotedblleft{}to end\textquotedblright{}? (\emph{\textbf{N-h-y}}, as in \emph{intahā}.) Which closer often carries relief? (\emph{\textbf{Akhīran}}, at last.)
\end{tcolorbox}
\section[\texorpdfstring{\ar{حكى} (ḥakā)}{حكى (ḥakā)}]{\ar{حكى} (ḥakā) — to tell a story}

\label{lesson:AR-C284-haka}



\begin{quote}
Say \textquotedblleft{}first\textquotedblright{}, \textquotedblleft{}then\textquotedblright{}, \textquotedblleft{}after that\textquotedblright{}. (\emph{Awwalan}, \emph{thumma}, \emph{baʿda dhālika}.) And the Arabic for a story. (\emph{Qiṣṣa}.)

You can put a story in order now. This is the verb for handing it on.
\end{quote}

\subsection*{You'll want to know: \ar{حكى}}
\textbf{\ar{حكى}} — \emph{ḥakā} — \textquotedblleft{}he told (a story)\textquotedblright{}, \textquotedblleft{}to tell, to relate\textquotedblright{}.

\begin{quote}
\textbf{\ar{حكت} \ar{جدتي} \ar{قصة} \ar{جميلة}.} — \emph{ḥakat jaddatī qiṣṣatan jamīla} — My grandmother told a lovely story.
\end{quote}

It ends like \emph{qaḍā}, and it bends the same way: \emph{ḥakaytu}, I told, as \emph{qaḍaytu}, I spent. The present is \emph{yaḥkī}, he tells, and with \emph{kāna} in front, \emph{kānat taḥkī}, she used to tell.

\emph{Qāla} reports what somebody said; \emph{ḥakā} passes on a whole story, a piece of news or something that happened.

\begin{cousinweb}
\textbf{\ar{حكى}} (\emph{ḥakā}) is built on the native Arabic root \emph{ḥ-k-y}, \textquotedblleft{}to relate, to narrate; to imitate\textquotedblright{} (Lane). Lane gives both senses side by side: to tell a story is to say again what was said or done, and to imitate is to do again what someone did.
\end{cousinweb}

\subsection*{Your turn}
Say these aloud:

\begin{itemize}
\raggedright
  \item \emph{ḥakā}, then \emph{ḥakaytu} and \emph{ḥakat}
  \item \textquotedblleft{}My grandmother told a lovely story\textquotedblright{} — \emph{ḥakat jaddatī qiṣṣatan jamīla}
  \item \textquotedblleft{}She used to tell a story every day\textquotedblright{} — \emph{kānat taḥkī qiṣṣatan kulla yawm}
  \item \textquotedblleft{}He said\textquotedblright{}, then \textquotedblleft{}he told a story\textquotedblright{} — \emph{qāla}, \emph{ḥakā qiṣṣa}
\end{itemize}

\begin{tcolorbox}[breakable,colback=teal!4,colframe=teal!35!black,title={Before you move on}]
What does \ar{حكى} mean? (\textquotedblleft{}To tell a story\textquotedblright{}.) \textquotedblleft{}I told\textquotedblright{}? (\emph{\textbf{Ḥakaytu}}, like \emph{qaḍaytu}.) \textquotedblleft{}He tells\textquotedblright{}? (\emph{\textbf{Yaḥkī.}}) What else can the root mean, in Lane? (\textbf{To imitate.})
\end{tcolorbox}
\section[\texorpdfstring{\ar{عندما} (ʿindamā)}{عندما (ʿindamā)}]{\ar{عندما} (ʿindamā) — when, in a story}

\label{lesson:AR-C284-indama}



\begin{quote}
Say \textquotedblleft{}she told a story\textquotedblright{}. (\emph{Ḥakat qiṣṣa}.) Then say \emph{ḥīnamā jāʾa, kataba}, \textquotedblleft{}when he came, he wrote\textquotedblright{}.

\emph{Ḥīnamā} has a twin that you will hear far more often in a story.
\end{quote}

\begin{grammarlens}[title={Grammar Lens: ʿindamā and the past}]
\textbf{\ar{عندما}} — \emph{ʿindamā} — \textquotedblleft{}when\textquotedblright{}. It is \emph{ʿinda}, \textquotedblleft{}at\textquotedblright{}, with \emph{mā} after it, as \emph{ḥīnamā} is \emph{ḥīna}, \textquotedblleft{}time\textquotedblright{}, with \emph{mā}. Put a past verb after it, and the clause names the moment the next event hangs on:

\begin{quote}
\textbf{\ar{عندما} \ar{وصلنا} \ar{إلى} \ar{المطار،} \ar{اتصلت} \ar{بزوجتي}.} — \emph{ʿindamā waṣalnā ilā al-maṭār, ittaṣaltu bi-zawjatī} — When we arrived at the airport, I phoned my wife.
\end{quote}

The \emph{ʿindamā} clause can follow as well: \emph{ittaṣaltu bi-zawjatī ʿindamā waṣalnā}. Either way, the arrival comes first and the call second.
\end{grammarlens}

\begin{grammarlens}[title={Grammar Lens: one event after another}]
\emph{Baʿda} means \textquotedblleft{}after\textquotedblright{} before a noun: \emph{baʿda al-ʿashāʾ}, after dinner. Before a whole clause it takes \emph{an}, \textquotedblleft{}that\textquotedblright{}, and the past:

\begin{quote}
\textbf{\ar{بعد} \ar{أن} \ar{أكلنا،} \ar{خرجنا}.} — \emph{baʿda an akalnā, kharajnā} — After we had eaten, we went out.
\end{quote}

So a story can be ordered two ways: by joiners, \emph{awwalan … thumma}, or by clauses, \emph{ʿindamā} and \emph{baʿda an}. Good stories mix them.
\end{grammarlens}

\subsection*{Your turn}
Say these aloud:

\begin{itemize}
\raggedright
  \item \emph{ʿindamā waṣalnā ilā al-maṭār, ittaṣaltu bi-zawjatī}
  \item the same two events the other way round — \emph{ittaṣaltu bi-zawjatī ʿindamā waṣalnā}
  \item \textquotedblleft{}After dinner\textquotedblright{}, then \textquotedblleft{}after we had eaten\textquotedblright{} — \emph{baʿda al-ʿashāʾ}, \emph{baʿda an akalnā}
  \item \emph{baʿda an akalnā, kharajnā, wa-fī an-nihāya rajaʿnā}
\end{itemize}

\begin{tcolorbox}[breakable,colback=teal!4,colframe=teal!35!black,title={Before you move on}]
Which two pieces make \emph{ʿindamā}? (\emph{\textbf{ʿInda}} and \emph{\textbf{mā}}.) Its twin? (\emph{\textbf{Ḥīnamā.}}) What does \emph{baʿda} take before a clause? (\emph{\textbf{An}}, then the past.) In \emph{ʿindamā waṣalnā, ittaṣaltu}, which came first? (\textbf{The arrival.})
\end{tcolorbox}
\section[(reading)]{(reading) — a trip to the sea, told in both pasts}

\label{lesson:AR-C284-qissat-rihla}



\begin{quote}
Chapter forty-two gave you a whole encounter to read: a greeting, a choice and a goodbye. This passage is a story, eighty-two words long, and it happened last summer.

Every word in it was taught before this lesson.
\end{quote}

\begin{tcolorbox}[breakable,skin=enhanced,colback=orange!6,colframe=orange!45!black,boxrule=0.5pt,arc=1mm,left=6pt,right=6pt,top=4pt,bottom=4pt,fonttitle=\bfseries,title={Read this}]
\begin{quote}
\ar{في} \ar{الصيف} \ar{الماضي} \ar{سافرت} \ar{مع} \ar{صديقي} \ar{إلى} \ar{قرية} \ar{صغيرة} \ar{قريبة} \ar{من} \ar{البحر}. \ar{كان} \ar{الطقس} \ar{جميلا}. \ar{عندما} \ar{وصلنا،} \ar{ذهبنا} \ar{أولا} \ar{إلى} \ar{الفندق}. \ar{ثم} \ar{أكلنا} \ar{السمك} \ar{في} \ar{مطعم} \ar{صغير،} \ar{وكان} \ar{الطعام} \ar{لذيذا}. \ar{وبعد} \ar{ذلك} \ar{سبحنا} \ar{في} \ar{البحر}.
\end{quote}

\begin{quote}
\ar{في} \ar{المساء} \ar{كنا} \ar{نشرب} \ar{الشاي،} \ar{وكان} \ar{صديقي} \ar{يحكي} \ar{قصة} \ar{طويلة}. \ar{كنت} \ar{أستمع} \ar{وأضحك}. \ar{ذات} \ar{مرة} \ar{كنا} \ar{نمشي} \ar{في} \ar{القرية} \ar{عندما} \ar{أمطرت}. \ar{بعد} \ar{أن} \ar{رجعنا} \ar{إلى} \ar{الفندق،} \ar{قرأنا} \ar{وشربنا} \ar{القهوة}.
\end{quote}

\begin{quote}
\ar{في} \ar{النهاية} \ar{رجعنا} \ar{بالقطار} \ar{إلى} \ar{المدينة}. \ar{وعندما} \ar{وصلت} \ar{إلى} \ar{المنزل،} \ar{حكيت} \ar{القصة} \ar{لجدتي}. \ar{كانت} \ar{الرحلة} \ar{قصيرة} \ar{وجميلة} \ar{جدا}.
\end{quote}

Read it once without stopping. Do not translate as you go; let the sentences arrive.

Then read it a second time and sort every verb into one of two piles: the scene, or what happened.
\end{tcolorbox}

\begin{culture}
The events: \emph{sāfartu}, \emph{waṣalnā}, \emph{dhahabnā}, \emph{akalnā}, \emph{sabaḥnā}, \emph{amṭarat}, \emph{rajaʿnā}, \emph{qaraʾnā}, \emph{ḥakaytu}. Each one happened and was over, and \emph{awwalan}, \emph{thumma}, \emph{baʿda dhālika} and \emph{fī an-nihāya} keep them in order.

The scene: the weather, the good food, the tea in the evening. \emph{Kāna} paints it, and so does \emph{kāna} with the present: \emph{kunnā nashrabu}, \emph{kāna ṣadīqī yaḥkī}, \emph{kuntu astamiʿu}. After \emph{kāna} an adjective takes \emph{-an}: \emph{jamīlan}, \emph{ladhīdhan}. Chapter two hundred and eighty-nine comes back to that.

In the rain, both pasts meet: \emph{kunnā namshī}, the walk still running, and \emph{ʿindamā amṭarat}, the rain that cut into it.
\end{culture}

\subsection*{Your turn}
\begin{itemize}
\raggedright
  \item \emph{Say it:} the passage aloud, once, without stopping
\end{itemize}

\begin{tcolorbox}[breakable,colback=teal!4,colframe=teal!35!black,title={Before you move on}]
How many words in that passage were new? (\textbf{None.})

Why is \emph{kunnā namshī} built with \emph{kāna}? (The walk \textbf{was still going on} when the rain came: it is scene.)

Who heard the story at the end? (\textbf{The grandmother}: \emph{ḥakaytu al-qiṣṣa li-jaddatī}.)
\end{tcolorbox}
\section[(review)]{Review — a story, in order}

\label{lesson:AR-R284-a-story-in-order}



\begin{quote}
Say \textquotedblleft{}I was reading\textquotedblright{}, then \textquotedblleft{}I was at the office\textquotedblright{}. (\emph{Kuntu aqraʾu}, \emph{kuntu fī al-maktab}.) And the Arabic for a page. (\emph{Ṣafḥa}.)
\end{quote}

\begin{grammarlens}[title={What you've built}]
Four lessons, and a story has its order. Joiners at the head of each step: \emph{awwalan}, \emph{thumma}, \emph{baʿda dhālika}, \emph{fī an-nihāya}. Clauses that name the moment: \emph{ʿindamā} with the past, and \emph{baʿda an} with the past. And \emph{ḥakā}, to hand the whole thing on.

One small story, in four steps:

\begin{quote}
\textbf{\ar{أولا} \ar{ذهبت} \ar{إلى} \ar{المكتب،} \ar{ثم} \ar{كتبت} \ar{رسالة}. \ar{عندما} \ar{رجعت،} \ar{حكيت} \ar{القصة} \ar{لصديقي}.} — \emph{awwalan dhahabtu ilā al-maktab, thumma katabtu risāla. ʿindamā rajaʿtu, ḥakaytu al-qiṣṣa li-ṣadīqī} — First I went to the office, then I wrote a letter. When I came back, I told my friend the story.
\end{quote}
\end{grammarlens}

\subsection*{Your turn}
Say these aloud:

\begin{itemize}
\raggedright
  \item \emph{awwalan dhahabtu ilā al-maktab, thumma katabtu risāla}
  \item \emph{ʿindamā rajaʿtu, ḥakaytu al-qiṣṣa li-ṣadīqī}
  \item \textquotedblleft{}After the client came, the pharmacist phoned\textquotedblright{} — \emph{baʿda an jāʾa al-ʿamīl, ittaṣala aṣ-ṣaydalī}
  \item the trip to the sea in three steps: where you went first, what you did then, how it ended
  \item one scene line for it — \emph{kāna aṭ-ṭaqs jamīlan}
\end{itemize}

\begin{tcolorbox}[breakable,colback=teal!4,colframe=teal!35!black,title={Before you move on}]
\textquotedblleft{}After that\textquotedblright{}? (\emph{\textbf{Baʿda dhālika.}}) \textquotedblleft{}When we arrived\textquotedblright{}, in a story? (\emph{\textbf{ʿIndamā waṣalnā.}}) \textquotedblleft{}After we had eaten\textquotedblright{}? (\emph{\textbf{Baʿda an akalnā.}}) \textquotedblleft{}I told the story\textquotedblright{}? (\emph{\textbf{Ḥakaytu al-qiṣṣa.}})
\end{tcolorbox}
